\subsection{Ant Colony Optimization (ACO)}

O algoritmo Ant Colony Optimization (ACO) é uma meta-heurística originada do Ant System, proposto por \citeonline{dorigo1996ant}, baseado no comportamento coletivo de formigas na busca por alimentos, principalmente no uso de informações por meio dos feromônios. Foi formulado como um método de otimização combinatória e inicialmente aplicado a problemas como o Traveling Salesman Problem (TSP). Posteriormente, diferentes algoritmos baseados no comportamento de colônias de formigas foram generalizados sob a estrutura da meta-heurística ACO \cite{dorigo1999ant}.

Fundamentalmente, o ACO funciona com base na construção incremental de soluções por um conjunto de formigas artificiais. Cada formiga seleciona probabilisticamente um dos componentes disponíveis para continuar a construção de uma solução. Essa escolha se baseia em duas informações principais: a trilha de feromônio $\tau$, que representa a experiência acumulada pelas formigas em decisões anteriores, e a informação heurística $\eta$, que indica a atratividade de uma escolha de acordo com as características do problema. A importância desses fatores é controlada pelos parâmetros $\alpha$ e $\beta$, sendo que as decisões com maiores valores de feromônio e de informação heurística têm maior tendência de serem selecionadas \cite{dorigo1999ant}. Após cada formiga realizar suas escolhas e finalizar a construção de sua solução, ocorre a atualização das trilhas de feromônio. Nesse processo, soluções de melhor qualidade podem receber um reforço na quantidade de feromônio, enquanto a evaporação reduz gradualmente os valores existentes, diminuindo a influência de decisões anteriores nas próximas iterações \cite{dorigo1999ant}.

No contexto da alocação de recursos e da virtualização de redes, abordagens baseadas em ACO têm sido investigadas como alternativas para problemas de otimização relacionados ao mapeamento e à alocação de recursos. \citeonline{fajjari2011vnACO} propuseram uma adaptação do ACO para o problema de Virtual Network Embedding (VNE), considerando a alocação de redes virtuais sobre uma infraestrutura física e avaliando a taxa de rejeição das requisições, o custo, a receita e a utilização média dos enlaces. Em outro contexto, \citeonline{wen2023power} utilizaram uma abordagem baseada em colônia de formigas para auxiliar na decisão de associação entre fluxos e slices, buscando maximizar o retorno de qualidade de serviço obtido por unidade de investimento em potência (ROI).

Considerando suas características, o ACO pode ser adaptado para solucionar problemas de fatiamento de redes tendo as possíveis alocações de recursos como soluções do processo de otimização. Nesse contexto, as formigas artificiais podem explorar diferentes configurações para atender às demandas existentes, enquanto os feromônios podem favorecer recursos presentes em soluções de melhor qualidade. Em ambientes de demanda elástica, esse processo pode considerar as alterações nas necessidades dos usuários ao longo do tempo, permitindo a busca por novas configurações de fatias conforme o estado dos recursos e as demandas apresentadas em cada período.